\printedpage{78}
% head: A. Grothendieck : Dérivateurs, Ch. XII, Édition des Universités de Montpellier II et Paris VII
\begin{enumerate}
\item[b)] I would like, for later applications, to have at my disposal a class of open immersions having the following properties:
\begin{enumerate}
\item[1°)] They are $W$-cofibrations, and
\item[2°)] If $U \overset{i}{\hookrightarrow} X$ is in the class, and if $\varepsilon \xrightarrow{j} I$ is the ``source-target'' inclusion of the discrete category $\{0,1\}$ into the interval
\[
\Phi_n = 0 \longrightarrow 1 \longleftarrow 2 \longrightarrow \dots 2n-1 \longleftarrow 2n
\]
$(j(0)=0$, $j(1)=2n)$, $n\geq 2$, the inclusion $\alpha$ of $(U\times I)\cup(X\times\varepsilon)$ into $X\times I$ is in the class.
\end{enumerate}
Now this last stability property is \textit{not} satisfied by Dwyer immersions. (Take $i : \{0\} \longrightarrow \boldsymbol{\Delta}^1 = \{0,1\}$, the inclusion $\alpha$ is not a Dwyer immersion.)
\end{enumerate}

\msp{page 115}

\subsection*{2.5\quad Complement: criterion for total asphericity of a $W$-proper category over another.}

Let
\[
f : X \longrightarrow I
\]
be a proper category over an index category $I$. If $I$ has a final object $i_0$, then $X/i_0$ is identified with $X$, and consequently the inclusion
\[
X_{i_0} \longrightarrow X
\]
of the fiber of $X$ at $i_0$ is $W$-aspheric \texttt{[rather }$W$\texttt{-coaspheric]}, and a fortiori is in $W$. Thus $X$ is aspheric if and only if $X_{i_0}$ is. We want to extend this result to a criterion for \textit{total} $W$-asphericity, more precisely to prove the

\textbf{Proposition 11:} \textit{Let $f : X \longrightarrow I$ be $W$-proper, and let $i_0$ be a final object of $I$. If $X_{i_0}$ is totally $W$-aspheric} (def. 1, page 88), \textit{the same is true of $X$.} (Provided that $W$ satisfies the connectedness axiom W(5).)

The proof will be given below, after some preliminaries of independent interest. First of all, let us recall that if $X$ is $W$-proper over $I$, then it gives rise to a functor
\[
\begin{aligned}
I \quad &\longrightarrow \quad \mathrm{Hot}_W\\
i \quad &\longmapsto \quad \mathrm{hot}_W(X_i)(\simeq \mathrm{hot}_W(X/i));
\end{aligned}
\]
composing it with the functor
\[
\pi_0 : \mathrm{Hot}_W \longrightarrow \mathrm{Ens}
\]
(provided that the connectedness axiom W(5) (p. 10) is satisfied)
% foot: Version 7 décembre 2002
\printedpage{79}
% head: A. Grothendieck : Dérivateurs, Ch. XII, Édition des Universités de Montpellier II et Paris VII
\msp{page 116}
we find a functor
\[
i \longmapsto \underbrace{\pi_0(X_i)}_{\simeq\, \pi_0(X/i)} : I \longrightarrow \mathrm{Ens}.
\]
This being so, we have the sorital result

\textbf{Proposition 12:} \textit{Let $f : X \longrightarrow I$ be $W$-proper, and suppose that $W$ satisfies the connectedness axiom} W(5) (otherwise, $W$-properness will be of no use to us \dots) (\footnote[81]{It is better to say: Let $f$ be $W_0$-proper, where $W_0$ is the localizer of 0-connectedness \dots}). \textit{We then have a canonical bijection}
\[
\pi_0(X) \xleftarrow{\ \sim\ } \varinjlim_{I} \pi_0(X_i), \tag*{$(*)$}
\]
\textit{where the arrow is induced by the canonical maps $\pi_0(X_i) \longrightarrow \pi_0(X)$.}

\textsc{Proof.} Passing to $f^\circ : X^\circ \longrightarrow I^\circ$, we are reduced to proving the same result for $X$ \textit{$W$-smooth} over $I$, with the only difference that the functor
\[
i \longmapsto \pi_0(X_i)
\]
on $I$ will be contravariant, i.e. a presheaf, which offers the advantage of the intuition attached to the formalism of presheaves. Let
\[
g : I \longrightarrow e, \qquad \text{whence } gf = p : X \longrightarrow e,
\]
it is immediate, by the standard calculation
\[
p_!(F) = \pi_0(X/F) \qquad (F \in \mathrm{Ob}\, X^\wedge)
\]
(where we have identified $e^\wedge$ with (Ens)), that we have
\[
\pi_0(X) = p_!(e_X) = g_!(f_!(e_X)).
\]
Now the presheaf
\[
i \longmapsto \pi_0(X_i) \qquad (\simeq \pi_0(i\backslash X), \text{ since } X\ W\text{-smooth over } I)
\]
\msp{page 117}
is none other than $f_!(e_X)$. On the other hand, for any presheaf $G$ on $I$, we have
\[
g_!(G) \simeq \pi_0(I/G) \simeq \varinjlim_{I} G(i),
\]
and applying this to $f_!(e_X)$, we find the formula $(*)$ of proposition 12.

Let us return to the situation with $f$ \textit{$W$-proper,} and consider a functor
\[
p : X \longrightarrow Y,
\]
we are interested in ``calculating''
\[
p_!(e_X) = (y \longmapsto \pi_0(y\backslash X)) \in \mathrm{Ob}\, Y^\wedge.
\]
Let us consider for this the following diagram
% foot: Version 7 décembre 2002
\printedpage{80}
% head: A. Grothendieck : Dérivateurs, Ch. XII, Édition des Universités de Montpellier II et Paris VII
\[
\begin{tikzcd}
y\backslash X_i \arrow[r] \arrow[d] & y\backslash X \arrow[r] \arrow[d] & y\backslash Y \arrow[d] \\
X_i \arrow[r] \arrow[d] & X \arrow[r, "p"] \arrow[d, "f"] & Y \\
\underset{\simeq\, e}{\{i\}} \arrow[r] & I\rlap{\quad ,} &
\end{tikzcd}
\]
where the three squares are cartesian. Since $y\backslash Y \longrightarrow Y$ is $W$-proper, the same is true of $y\backslash X \longrightarrow X$ by base change, hence of the composite $y\backslash X \longrightarrow X \longrightarrow I$ by transitivity. The diagram shows on the other hand that the fiber of $y\backslash X$ at $i \in \mathrm{Ob}\, I$ is identified with $y\backslash X_i$.

\msp{page 118}
Proposition 12 therefore gives us the formula
\[
\pi_0(y\backslash X) = \varinjlim_{I} \pi_0(y\backslash X_i),
\]
which formula cannot help being functorial in $y$, just as that of prop. 12 is in $\underline{\mathrm{Prop}}_W I$. Now for $y$ variable in $Y$, the presheaf on $Y$
\[
y \longmapsto \pi_0(y\backslash X_i)
\]
is none other than $p_{i\,!}(e_{X_i})$, where
\[
p_i : X_i \longrightarrow Y
\]
is the composite $X_i \longrightarrow X \xrightarrow{p} Y$ (induced by $p$). We have thus proved

\textbf{Corollary 1} (to prop. 12). \textit{Let $f : X \longrightarrow I$ be $W$-proper} (\footnote[82]{say $W_0$-proper \dots}), \textit{and $p : X \longrightarrow Y$ any functor. Let, for every $i \in \mathrm{Ob}\, I$, $p_i : X_i \longrightarrow Y$ be the functor induced on the fiber $X_i$ of $X$ at $i$. We then have a natural} functorial \textit{law}
\[
i \longmapsto p_{i\,!}(e_{X_i}) : I \longrightarrow Y^\wedge,
\]
\textit{and we have the formula}
\[
p_!(e_X) \xleftarrow{\ \sim\ } \varinjlim_{I} p_{i\,!}(e_{X_i}), \tag*{$(*)$}
\]
\textit{where the arrow is induced by the canonical morphisms}
\[
p_{i\,!}(e_{X_i}) \longrightarrow p_!(e_X)
\]
% foot: Version 7 décembre 2002
\printedpage{81}
% head: A. Grothendieck : Dérivateurs, Ch. XII, Édition des Universités de Montpellier II et Paris VII
\msp{page 119}
\textit{induced by the inclusions $X_i \longrightarrow X$.}

The functorial law of $p_{i\,!}(e_{X_i})$ in $i$ has been made explicit in the course of the proof.

Let us note that when $X$ over $I$ is cofibered and split, i.e. comes from a functor
\[
i \longmapsto X_i, \quad I \longrightarrow \mathrm{Cat},
\]
then the functorial law of $p_{i\,!}(e_{X_i})$ in $i$ is made explicit in a particularly pleasant way: to an arrow
\[
\alpha : i \longrightarrow j, \qquad \text{whence } \alpha_* : X_i \longrightarrow X_j,
\]
is associated
\[
p_!(\alpha) : p_{i\,!}(e_{X_i}) \longrightarrow p_{j\,!}(e_{X_j}),
\]
noting that we have a diagram
\[
\begin{tikzcd}
X_i \arrow[rr, "\alpha_*"] \arrow[dr, "p_i"', ""{name=A}] & & X_j \arrow[dl, "p_j", ""'{name=B}] \\
& Y &
\arrow[from=A, to=B, bend left, "\lambda_\alpha"] % [resolved: see DER_XIIb-a_corrections.csv]
\end{tikzcd}
\]
with a commutativity arrow
\[
\lambda_\alpha : p_i \longrightarrow p_j\alpha_*
\]
(not \texttt{[an]} isomorphism in general), whence
\[
p_{i\,!}(e_{X_i}) \xrightarrow{(\lambda_\alpha)_!} p_{j\,!}\underbrace{\alpha_{*\,!}(e_{X_i})}_{\longrightarrow\, e_{X_j}} \longrightarrow p_{j\,!}(e_{X_j}),
\]
no more doubt that's the one! (\footnote[83]{This arrow \texttt{[the arrow }$(\lambda_\alpha)_!$\texttt{]} should be made explicit with care!}).

\msp{page 120}
\textbf{Corollary 2:} \textit{Under the conditions of corollary 1, suppose that $I$ has a final object $i_0$. Then}
\[
(p_{i_0})_!(e_{X_{i_0}}) \longrightarrow p_!(e_X)
\]
\textit{is an isomorphism in $Y^\wedge$.}

Whence at once proposition 11: if $X_{i_0}$ is totally $W$-aspheric, then $X$ is too, since for every $p : X \longrightarrow Y$, $p_!(e_X) \simeq p_{i_0\,!}(X_{i_0})$, and this latter element of $Y^\wedge$ is $W$-aspheric % [sic] p_{i_0 !}(X_{i_0}) -> p_{i_0 !}(e_{X_{i_0}}), elsewhere the argument of p_{i!} is e_{X_i}
by the hypothesis on $X_{i_0}$.

Let us note the special case:

\textbf{Corollary 3.} \textit{Let $Z$ be a totally $W$-aspheric object in} Cat. \textit{Then its product $Z\times I$ with any $I$ in} Cat \textit{having a final object is also totally $W$-aspheric.}
% foot: Version 7 décembre 2002
\printedpage{82}
% head: A. Grothendieck : Dérivateurs, Ch. XII, Édition des Universités de Montpellier II et Paris VII

\msp{page 121}

\section*{3\quad Amalgamated sums and integration\quad \begin{tikzcd}[ampersand replacement=\&, row sep=small, column sep=small, baseline=(\tikzcdmatrixname-2-1.base)]
\& X_1 \\
\int X_0 \arrow[ur] \arrow[dr] \& \\
\& X_2
\end{tikzcd}}

\subsection*{3.1\quad Canonical reduction to the case of a diagram of closed immersions}

Let us consider a diagram
\[
\begin{tikzcd}[row sep=small, column sep=small]
& X_1 \\
X_0 \arrow[ur, "\alpha_1"] \arrow[dr, "\alpha_2"'] & \\
& X_2
\end{tikzcd}
\tag*{$(*)$}
\]
in Cat. We can regard it as a contrafunctor
\[
\Phi^\circ \longrightarrow \mathrm{Cat},
\]
where
\[
\Phi = \left( \vcenter{\hbox{\begin{tikzcd}[row sep=small, column sep=small]
& 1 \arrow[dl] \\
0 & \\
& 2 \arrow[ul]
\end{tikzcd}}} \right),
\]
thus as defining a (split) fibered category $\mathscr{X}$ over $\Phi$. It is sometimes convenient to represent this category $\mathscr{X}$ by the ``integration sign'' % [sic] integration -> intégration, accent missing in the print
\[
\begin{tikzcd}[row sep=small, column sep=small]
& X_1 \\
\int X_0 \arrow[ur] \arrow[dr] & \\
& X_2,
\end{tikzcd}
\]
and regard it as a sort of ``homotopy amalgamated sum'' of $X_1$ and $X_2$ under $X_0$, or as a ``homotopy direct limit'' of the original diagram $(*)$.

\msp{page 122}
By this one is to understand that (for a given $W \subseteq \mathrm{Fl}(\mathrm{Cat})$) the image $\mathrm{hot}_W(\mathscr{X})$ of $\mathscr{X}$ in $\mathrm{Hot}_W$ has a remarkable ``homotopy meaning'' in terms of $(*)$, more precisely, the diagram $(*)$ defines an object $\xi$ of $\mathrm{HOT}_W(\Phi)$ (which is a category of fractions of the category of all diagrams of type $(*)$ \dots), and if
\[
\pi : \Phi \longrightarrow e
\]
denotes the absolute projection of $\Phi$, we have
\[
\pi_!(\xi) = \mathrm{hot}_W(\mathscr{X}),
\]
% foot: Version 7 décembre 2002